\documentclass[11pt]{article}
\usepackage{amsmath,amssymb,amsthm,mathtools}
\usepackage[margin=1in]{geometry}

\newtheorem{theorem}{Theorem}
\newtheorem{definition}{Definition}
\newtheorem{lemma}{Lemma}

\title{Human-Readable Proof Outline for the Kasami APN Theorem (Odd Dimension)}
\author{RequestProject / Kasami APN}
\date{}

\begin{document}
\maketitle

\section{What is proved}
The Lean theorem \texttt{kasami\_is\_apn} proves that the Kasami power map
\[
  f(x) = x^{d(k)}, \qquad d(k)=2^{2k}-2^k+1,
\]
is APN on a field $F$ of characteristic $2$ under the assumptions:
\begin{itemize}
  \item $|F|=2^n$,
  \item $n$ odd,
  \item $1<k<n$,
  \item $k$ odd,
  \item $\gcd(k,n)=1$.
\end{itemize}

Among these assumptions, the odd-dimension condition ($n$ odd) becomes essential at Layer 5,
where it is used to obtain the coprimality input needed to construct the linking exponent $e'$. 

Also, at Layer 6, the Gold-map bijectivity is used to cancel a $(2^k+1)$-power, which is only valid under the coprimality assumption $\gcd(2^k+1,2^n-1)=1$.

\section{APN meaning and reduction strategy}
For a function $f:F\to F$, APN means that for every nonzero $a$ and every collision
\[
  f(x+a)+f(x)=f(y+a)+f(y),
\]
the only possible outputs are $y=x$ or $y=x+a$.

The proof does not attack this directly for all $a$. It first uses the normalization lemma
\texttt{apn\_of\_normalized}: it is enough to prove the same uniqueness statement for the
normalized equation with $a=1$:
\[
  (x+1)^d+x^d=(y+1)^d+y^d \implies y=x \text{ or } y=x+1.
\]

So the whole theorem is reduced to controlling collisions of the normalized differential.

\section{Layered proof architecture}
The Lean file is organized in layers. The theorem-level dependency chain is:
\[
\begin{array}{c}
\texttt{kasami\_is\_apn} \\
\uparrow \\
\texttt{apn\_of\_normalized} \\
\uparrow \\
\texttt{kasami\_collision\_forces\_equal\_u} \\
\uparrow \\
\texttt{phi\_injective\_on\_units} \\
\uparrow \\
\texttt{exists\_linking\_exp}
\end{array}
\]

Additional inputs are the Artin--Schreier identity, the key Kasami polynomial identity,
Gold-map bijectivity, and a quadratic-kernel fact in characteristic $2$.

\subsection*{Layer 1: Artin--Schreier trace identity}
\texttt{truncTrace\_artin\_schreier} proves
\[
  L_k(x^2+x)=x^{2^k}+x.
\]
This is the basic conversion between truncated trace expressions and explicit powers.

\subsection*{Layer 2: Key Kasami identity}
\texttt{kasami\_key\_identity} proves
\[
\bigl((x+1)^d+x^d+1\bigr)(x^2+x)^{2^k}=(x^{2^k}+x)^{2^k+1}.
\]
This is the bridge from differential terms to a form suited for the $\Phi$-map analysis.

\subsection*{Layer 3: Gold coprimality and bijectivity}
\texttt{gold\_coprime} gives
\[
  \gcd(2^k+1,2^n-1)=1,
\]
and \texttt{gold\_pow\_bijective} deduces that
\[
  y\mapsto y^{2^k+1}
\]
is bijective on $F$.

This bijectivity is later used to cancel an outer $(2^k+1)$-power.

\subsection*{Layer 4: Arithmetic transfer identity}
\texttt{kasami\_arith\_identity} is a modular arithmetic identity tailored to the exponent transfer step.
It aligns the Kasami exponent arithmetic with the Dempwolff--Mueller bijection hypothesis.

\subsection*{Layer 5: Existence of a linking exponent $e'$}
\texttt{exists\_linking\_exp} constructs $e'$ with two modular constraints:
\[
e'(2^k+1)\equiv 2^n-1-2^k \pmod{2^n-1},
\]
\[
(2^{n-1}-2^{k-1}-1)e'\equiv 2^{k-1} \pmod{2^n-1}.
\]
Interpretation: $e'$ is the arithmetic compatibility parameter that lets the map
\[
\Phi(u)=\frac{L_k(u)^{2^k+1}}{u^{2^k}}
\]
be rewritten as a composition of two bijective-style maps.

This is the first point in the proof where the hypothesis ``$n$ odd'' is genuinely needed:
it enters through the Gold coprimality step used in constructing $e'$.

\subsection*{Layer 6: Injectivity of $\Phi$ on nonzero elements}
\texttt{phi\_injective\_on\_units} proves that for $u,v\neq 0$, a collision in the $\Phi$-equation forces $u=v$.
The proof combines:
\begin{itemize}
  \item the exponent link from Layer 5,
  \item Gold-map bijectivity to cancel $(2^k+1)$ powers,
  \item Dempwolff--Mueller theorem \texttt{LxXk'\_bijective} to conclude equality of inputs.
\end{itemize}

The Gold-map bijectivity is only valid under the coprimality assumption $\gcd(2^k+1,2^n-1)=1$.

\subsection*{Layer 7: Differential collision implies equal Artin--Schreier values}
\texttt{kasami\_collision\_forces\_equal\_u} proves:
\[
  (x+1)^d+x^d=(y+1)^d+y^d \implies x^2+x=y^2+y.
\]
It applies Layer 2 to rewrite the collision, then:
\begin{itemize}
  \item handles zero cases directly,
  \item uses Layer 6 injectivity in the nonzero case,
  \item uses Layer 1 for trace-power conversion.
\end{itemize}

This is the key bridge from normalized differential collisions to a quadratic relation in $x+y$.

\subsection*{Layer 8: Characteristic-2 kernel and normalization lemma}
There are two lemmas used in the final step:
\begin{itemize}
  \item \texttt{sq\_add\_self\_eq\_zero\_char2}:
  \[
    u^2+u=0 \iff u\in\{0,1\}.
  \]
  \item \texttt{apn\_of\_normalized}: proving the normalized differential uniqueness
  implies full APN.
\end{itemize}

\subsection*{Layer 9: Final theorem}
In \texttt{kasami\_is\_apn}, the final argument is short:
\begin{enumerate}
  \item Apply \texttt{apn\_of\_normalized}; reduce to the $a=1$ collision equation.
  \item Use Layer 7 to get $x^2+x=y^2+y$.
  \item Rearrange to
  \[
    (x+y)^2+(x+y)=0.
  \]
  \item Use Layer 8 kernel lemma to conclude $x+y=0$ or $x+y=1$.
  \item Therefore $y=x$ or $y=x+1$, exactly the APN condition.
\end{enumerate}

\section{One-paragraph summary}
The proof is modular. It first reduces APN to a normalized differential equation, then converts that equation into a trace/exponent form via the Kasami identity, proves injectivity of the associated $\Phi$-map on nonzero elements using a specially constructed linking exponent and known bijectivity theorems, and finally translates the resulting equality into the two-point alternative $y\in\{x,x+1\}$ through a characteristic-2 kernel identity. This closes the APN theorem for Kasami exponents in odd dimension under the standard coprimality assumptions.

\end{document}